% !TEX program = pdflatex
\documentclass[conference]{IEEEtran}

\usepackage{hyperref}
\usepackage{amsmath,amssymb,amsfonts}
\usepackage{graphicx}
\usepackage{booktabs}
\usepackage{multirow}
\usepackage{array}
\usepackage{url}
\usepackage{cite}
\usepackage{float}
\usepackage{setspace}

\title{
A Data-Driven Approach to Customer Lifetime Value Prediction and Marketing Optimization Using Machine Learning and A/B Testing
\thanks{This work was conducted as part of the DDM -- Data-Driven Marketing course.}
}

\author{
\IEEEauthorblockN{
\begin{tabular}{ll}
Group 2: & Nguyen Thi Van Anh, Vu Bui Dinh Tung, \\
         & Dao Khanh Linh, Nguyen Hong Nhung, Thieu Dieu Thuy
\end{tabular}
}

\IEEEauthorblockA{
DSEB-MFE, NEU \\
Course: DDM -- Data-Driven Marketing \\
}
}

\begin{document}
\maketitle

\begin{abstract}
Modern e-commerce platforms generate massive clickstream data characterized by sparsity, behavioral heterogeneity, and highly imbalanced transaction distributions. Traditional revenue prediction models often struggle to capture the complex relationship between customer engagement and purchasing behavior, particularly when only a small proportion of sessions generate revenue. This study proposes a business-oriented machine learning framework for customer conversion and revenue prediction using the Google Merchandise Store dataset from Google Analytics.

The proposed pipeline combines preprocessing of semi-structured JSON logs, exploratory behavioral analysis, and advanced feature engineering techniques including Recency-Frequency-Monetary (RFM) metrics, session velocity indicators, and temporal browsing patterns. To address the zero-inflated nature of transaction revenue, the study implements a two-stage XGBoost hurdle architecture consisting of a calibrated classifier for purchase propensity estimation and a weighted regressor for conditional revenue prediction.

Unlike conventional approaches evaluated only through statistical accuracy, this framework incorporates business-centric metrics such as Revenue Capture Curves and Lift analysis to measure commercial effectiveness. In addition, the predictive framework is connected to causal validation strategies through A/B testing, enabling practical deployment for marketing optimization and customer targeting.

Experimental analysis demonstrates that behavioral engagement, acquisition channels, and temporal interaction patterns provide strong predictive signals for both conversion probability and revenue generation. The resulting framework offers an interpretable and scalable solution for intelligent e-commerce analytics and customer value prediction.
\end{abstract}

\begin{IEEEkeywords}
E-commerce analytics, clickstream analysis, XGBoost, hurdle models, customer conversion prediction, revenue forecasting, feature engineering, behavioral analytics.
\end{IEEEkeywords}

%--------------------------------------------------
\section{Introduction}

The rapid growth of e-commerce platforms has transformed customer interactions into large-scale behavioral data ecosystems. Every browsing session, product interaction, and marketing touchpoint generates valuable clickstream information that can potentially reveal purchasing intent and customer value. However, extracting meaningful insights from such data remains challenging due to semi-structured formats, noisy traffic behavior, and severe transaction imbalance.

The Google Merchandise Store dataset provides a realistic analytics environment collected directly from Google Analytics, containing session-level information related to devices, traffic sources, geographic attributes, and transaction activities. A major challenge of the dataset is that only a very small proportion of sessions produce revenue, while the majority represent non-purchasing browsing behavior. This imbalance makes conventional single-stage regression models ineffective for commercial prediction tasks.

To address these limitations, this study proposes a framework built upon four key components. First, advanced feature engineering techniques are applied to extract behavioral signals from clickstream data, including RFM indicators, session velocity, engagement depth, and temporal browsing patterns. Second, a two-stage hurdle architecture based on XGBoost is introduced to separately model purchase probability and conditional transaction value through a calibrated classifier and weighted regressor. Third, the framework adopts business-centric evaluation metrics such as Revenue Capture Curves and Lift analysis to measure practical commercial performance beyond traditional statistical accuracy. Finally, the study connects predictive modeling with causal validation through A/B testing strategies to support real-world marketing optimization and customer targeting decisions.

The objective of this research is to develop an interpretable and scalable analytical framework capable of transforming raw behavioral analytics data into actionable business intelligence for conversion prediction and revenue optimization.

%--------------------------------------------------
\section{Exploratory Data Analysis}

The exploratory data analysis (EDA) process was conducted to understand the structural characteristics of the Google Merchandise Store dataset and uncover behavioral patterns associated with customer engagement, traffic acquisition, and transaction generation. Since the dataset originates from Google Analytics logs, the analysis focused not only on conventional statistical exploration but also on interpreting user behavior within a real-world e-commerce environment.

Unlike clean benchmark datasets commonly used in academic experiments, the Google Merchandise Store dataset contains highly heterogeneous information embedded within nested JSON structures. Therefore, the initial stage of EDA concentrated on understanding the overall architecture of the data before performing any modeling-related operations.

\subsection{Dataset Structure and Semi-Structured Characteristics}

The dataset consists of session-level records collected from user interactions on the Google Merchandise Store platform. Each observation represents a single visit session and contains behavioral, geographical, transactional, and acquisition-related information.

A major characteristic of the dataset is the presence of nested JSON attributes such as \textit{device}, \textit{geoNetwork}, \textit{totals}, and \textit{trafficSource}. These columns contain hierarchical structures rather than standard tabular values, requiring additional parsing and flattening procedures before analysis.

From an analytical perspective, this structure reflects the complexity of modern digital analytics systems where customer behavior is distributed across multiple contextual dimensions including device usage, browsing source, geographic origin, and engagement intensity.

After flattening the JSON columns, the dataset revealed several major groups of variables:

\begin{itemize}
    \item Device-related information such as browser type, operating system, and device category.
    \item Geographic indicators including country, city, and continent.
    \item Engagement metrics such as hits, pageviews, and visit counts.
    \item Traffic acquisition variables including referral source, campaign, and channel grouping.
    \item Transaction-related variables associated with purchases and revenue generation.
\end{itemize}

This structure provides a comprehensive representation of the customer journey from acquisition to conversion.

\subsection{Data Quality Inspection and Missing Value Analysis}

The next stage of EDA focused on evaluating data quality and identifying inconsistencies within the raw analytics logs.

The analysis revealed that multiple attributes contained substantial proportions of missing values or placeholder categories such as \texttt{(not set)} and \texttt{(not provided)}. Missingness was particularly common in variables associated with marketing campaigns and geographic granularity.

However, further inspection indicated that these missing values were not entirely random. For example, direct traffic sessions frequently lacked campaign metadata because users entered the website manually rather than through tracked advertising campaigns. Similarly, several geographic fields were unavailable due to browser privacy restrictions or incomplete tracking information.

This observation is important because it demonstrates that missing values in behavioral analytics datasets often carry implicit business meaning. Consequently, the analysis treated missingness as part of the behavioral context rather than simply removing incomplete observations.

In addition to missing values, the dataset also contained highly sparse variables and several low-variance features with nearly constant values across sessions. Such attributes contribute minimal analytical value and increase dimensional complexity.

The EDA process therefore identified candidate variables for removal while preserving behaviorally informative attributes.

\subsection{Revenue Distribution and Transaction Imbalance}

One of the most significant findings obtained during exploration was the extremely imbalanced nature of transaction activity.

The variable \textit{totals\_transactionRevenue} contained non-zero values in only approximately 1.08\% of all sessions. This indicates that the vast majority of users visited the platform without completing a purchase.

From a business perspective, this pattern is consistent with real-world e-commerce ecosystems where browsing behavior heavily outweighs actual conversion activity. Most visitors explore products, compare options, or navigate informational pages before leaving the platform.

From a data science perspective, this imbalance creates substantial modeling challenges because conventional algorithms tend to prioritize the dominant zero-revenue observations while underrepresenting high-value purchasing sessions.

Furthermore, the revenue distribution exhibited strong positive skewness, where a small subset of customers generated disproportionately large transaction values.

To better understand this behavior, the analysis examined both the raw revenue distribution and the logarithmically transformed revenue variable:

\begin{equation}
 y' = \log(1 + revenue)
\end{equation}

The transformed distribution significantly reduced the influence of extreme outliers and revealed clearer patterns among revenue-generating sessions.

\subsection{User Engagement and Browsing Behavior}

After analyzing the target variable, the EDA process investigated customer engagement patterns through variables such as \textit{totals\_hits}, \textit{totals\_pageviews}, and session frequency.

The distributions of hits and pageviews were heavily right-skewed. Most users interacted with only a small number of pages, while a limited subset demonstrated very deep browsing activity.

A detailed comparison between purchasing and non-purchasing sessions revealed that high-engagement users consistently generated greater conversion.

%--------------------------------------------------
\section{Modeling}

\subsection{Feature Engineering}
After completing the exploratory analysis stage, the next objective was to transform raw session-level clickstream data into customer-level behavioral representations suitable for downstream machine learning models. Since the Google Analytics dataset is highly sparse and dominated by short non-converting sessions, the feature engineering pipeline focused on aggregating historical user behavior across time windows rather than modeling isolated sessions independently.

To avoid temporal leakage, all features were generated inside rolling time folds consisting of a historical feature window and a future target window. For each fold, customer behavioral features were constructed exclusively from past sessions before predicting future revenue activity.

\subsubsection{Categorical Standardization and Target Encoding}

Several categorical variables extracted from the nested JSON structures contained extremely high cardinality, including geographic attributes, traffic acquisition sources, browser information, and campaign identifiers.

Instead of applying direct one-hot encoding, the pipeline first performed Top-$N$ grouping using only categories observed frequently in the training sessions. Rare categories were grouped into a common \texttt{Other} class in order to reduce sparsity and stabilize feature distributions across folds.

The standardized categorical variables included:

\begin{itemize}
    \item \textit{geoNetwork\_city}
    \item \textit{geoNetwork\_country}
    \item \textit{geoNetwork\_region}
    \item \textit{geoNetwork\_metro}
    \item \textit{geoNetwork\_subContinent}
    \item \textit{trafficSource\_source}
    \item \textit{trafficSource\_campaign}
    \item \textit{trafficSource\_medium}
    \item \textit{channelGrouping}
    \item \textit{device\_browser}
    \item \textit{device\_operatingSystem}
    \item \textit{device\_deviceCategory}
\end{itemize}

After grouping, target encoding was applied at the session level using the historical mean transaction revenue learned from training sessions only. This approach preserves behavioral information while avoiding the dimensional explosion caused by sparse categorical encoding.

\subsubsection{Customer-Level Aggregation Features}

Following categorical preprocessing, session-level interactions were aggregated using the \textit{fullVisitorId} identifier to construct long-term behavioral summaries for each customer.

Multiple aggregation operators including \texttt{sum}, \texttt{mean}, \texttt{max}, and \texttt{count} were applied to engagement and transaction variables such as:

\begin{itemize}
    \item \textit{totals\_hits}
    \item \textit{totals\_pageviews}
    \item \textit{totals\_timeOnSite}
    \item \textit{visitNumber}
    \item \textit{totals\_transactions}
    \item \textit{totals\_transactionRevenue}
\end{itemize}

This aggregation process generated behavioral indicators representing cumulative engagement intensity, browsing consistency, and historical transaction activity across multiple sessions.

\subsubsection{RFM and Customer Lifetime Features}

To capture long-term purchasing behavior, the framework additionally constructed Recency-Frequency-Monetary (RFM)-style features.

Recency variables measured the temporal distance between customer visits, while frequency indicators captured repeated interaction activity across the feature window. Monetary behavior was represented through aggregated historical revenue statistics such as total, average, and maximum transaction revenue.

Beyond standard RFM metrics, probabilistic customer lifetime features were generated using the BG/NBD and Gamma-Gamma models from the \texttt{lifetimes} library. These models produced forward-looking behavioral estimates including:

\begin{itemize}
    \item probability of customer survival (\textit{p\_alive}),
    \item expected future purchases (\textit{expected\_purchases\_62d}),
    \item expected average monetary value (\textit{expected\_monetary}).
\end{itemize}

These features provide a probabilistic representation of future customer value rather than relying solely on historical summaries.

\subsubsection{Session Velocity and Behavioral Intensity}

The pipeline also engineered velocity-based behavioral indicators to capture interaction intensity within browsing sessions.

The constructed variables included:

\begin{itemize}
    \item \textit{vel\_pageviews}
    \item \textit{vel\_hits}
    \item \textit{vel\_sessions}
    \item \textit{funnel\_depth}
\end{itemize}

These variables describe how actively customers interact with the platform over time. High-conversion users were frequently associated with deeper funnel traversal and more concentrated browsing activity compared to short exploratory sessions.

\subsubsection{Temporal Behavioral Features}

Additional temporal features were extracted to capture cyclical user activity patterns.

The engineered variables included:

\begin{itemize}
    \item \textit{pct\_weekend}
    \item \textit{pct\_evening}
    \item \textit{recency\_days}
    \item \textit{lifetime\_days}
\end{itemize}

These features represent customer browsing preferences across different time periods and allow the model to learn temporal variations in purchasing behavior.

% \subsubsection{Feature Stabilization and Outlier Handling}

% Finally, the engineered numerical features were stabilized using clipping and transformation procedures.

% Because engagement and monetary variables exhibited extreme heavy-tailed distributions, several features were clipped at the 99th percentile using thresholds calculated exclusively from training folds. The clipping process was applied to engagement variables, transaction statistics, velocity indicators, and monetary features such as \textit{expected\_monetary}.

% This step reduced outlier dominance while preserving the relative ordering of high-value customers, resulting in a more stable feature space for the downstream two-stage XGBoost framework.


\end{document}